\SourceLeaf{128}
\begin{center}
\textbf{ANZEIGE.}
\end{center}

\medskip
\noindent\rule{\linewidth}{.3pt}

\noindent Allgemeine Literatur-Zeitung vom Jahre 1808. Halle-Leipzig 1808. Nr.\,45. Februar 12. S.\,353--358.

\noindent\rule{\linewidth}{.3pt}
\medskip

Dresden, in der Waltherschen Hofbuchhandlung: \textsc{Leonellis} \emph{logarithmische Supplemente, als ein Beitrag, Mängel der gewöhnlichen Logarithmentafeln zu ersetzen. Aus dem Französischen nebst einigen Zusätzen von} \textsc{Gottfried Wilhelm Leonhardi}, \emph{Souslieutenant beim kurfürstl. sächsischen Feldartilleriecorps.} 1806. 88 S. in Octav.

Diese kleine Schrift enthält zwei von einander unabhängige Abhandlungen, die indessen in so fern einen gemeinschaftlichen Zweck haben, als beide etwas zu leisten bestimmt sind, was sich mit den gewöhnlichen logarithmischen Tafeln ohne anderweitige Hülfsmittel entweder nicht so vollkommen, oder nicht so bequem erreichen lässt. Die eine soll nemlich dazu dienen, vermittelst einiger sehr geschmeidiger, hier zugleich mitgelieferter, Hülfstafeln, zu jeder gegebenen Zahl ohne zu grosse Mühe den Logarithmen auf mehrere Decimalen, als die gewöhnlichen Tafeln verstatten, zu berechnen. Die andere entwickelt die Idee einer besondern Tafel, vermittelst welcher die Logarithmen von Summen oder Differenzen zweier bloss durch ihre Logarithmen gegebenen Grössen durch eine einzige Operation mit einer Bequemlichkeit sollen bestimmt werden können, wie sie bei andern Verfahrungsarten nicht Statt findet; und zwar, des Vfs. Plane nach, gleichfalls mit einer ungewöhnlich grossen Anzahl von Decimalen (14); von dieser Tafel ist jedoch nur erst eine Probe beigefügt. Da das Numerische der Logarithmen und jede sich darauf beziehende Erleichterung jedem, der viel mit Zahlenrechnungen zu thun hat, von grosser Wichtigkeit ist: so wird es sich wohl der Mühe verlohnen, diesen Untersuchungen eine nicht
